\documentclass[11pt]{article}
\usepackage[T1]{fontenc}
\usepackage{textcomp}
\usepackage[margin=0.75in]{geometry}
\usepackage{amsmath,amssymb,amsthm}
\usepackage{array,booktabs}
\usepackage{hyperref}
\setlength{\parindent}{0pt}
\newtheorem{theorem}{Theorem}
\theoremstyle{definition}
\newtheorem{definition}{Definition}
\title{Flow Proof Artifact: prop-44-apply-parallelogram-to-line}
\author{Generated by \texttt{flow doc proof}}
\date{Flow Proof Book}
\begin{document}
\maketitle
To a given straight line to apply a parallelogram equal to a given triangle in a given angle.\\[0.5em]
\textbf{Source.} euclid — Elements, Book I, Proposition 44\\[0.5em]
\bigskip
\noindent{\large \textbf{Derived fact 1.} \textit{To a given straight line to apply a parallelogram equal to a given triangle in a given angle} \hfill the Euclidean plane \cdot Euclid Book I \cdot Proposition 44: apply a parallelogram equal to a triangle to a line}\par
\smallskip\noindent\textit{Coordinate: the Euclidean plane \cdot Euclid Book I \cdot Proposition 44: apply a parallelogram equal to a triangle to a line}\par
\smallskip\noindent\textit{Source: euclid — Elements, Book I, Proposition 44}\par
\smallskip\noindent\textit{Built on: proposition 23: copy a given angle onto a line at a point, for Book I of the Elements in on the Euclidean plane, proposition 42: a parallelogram equals a triangle on an equal base, for Book I of the Elements in on the Euclidean plane}\par
\medskip
\noindent\textbf{Goal.} To a given straight line to apply a parallelogram equal to a given triangle in a given angle\par
\noindent$\displaystyle \text{parallelogram on AB equals triangle C}$\par
\medskip
\noindent
\begin{tabular}{@{} >{\raggedright\arraybackslash}p{0.06\textwidth} >{\raggedright\arraybackslash}p{0.47\textwidth} >{\raggedleft\arraybackslash}p{0.41\textwidth} @{}}
\textbf{\#} & \textbf{Proof} & \textbf{Mathematics} \\
\midrule
\textbf{1.} & We prove that proposition 44: apply a parallelogram equal to a triangle to a line for Book I of the Elements in the Euclidean plane. & \\[0.45em]
\textbf{2.} & Let angle\_at\_A = copied\_angle\_DCE\_at\_point\_A. & \\[0.45em]
\textbf{3.} & Let parallelogram\_on\_AB = parallelogram\_with\_angle\_at\_A. & \\[0.45em]
\textbf{4.} & We invoke the derived fact: segment\_AB\_given\_and\_triangle\_C\_and\_angle\_DCE. & \\[0.45em]
\textbf{5.} & We invoke the derived fact governing Book I of the Elements in the Euclidean plane: proposition 23: copy a given angle onto a line at a point, for Book I of the Elements in on the Euclidean plane. & \\[0.45em]
\textbf{6.} & We invoke the derived fact governing Book I of the Elements in the Euclidean plane: proposition 42: a parallelogram equals a triangle on an equal base, for Book I of the Elements in on the Euclidean plane. & \\[0.45em]
\textbf{7.} & From step 2, step 3, step 4, step 5, and step 6, by the chain of results established in Book I (step 2, step 3, step 4, step 5, and step 6), we obtain parallelogram on AB equals triangle C. Hence proven. & $\displaystyle \text{parallelogram on AB equals triangle C}$ \\[0.45em]
\bottomrule
\end{tabular}
\label{thm:Geometry-Euclid Book I-proposition_44_apply_a_parallelogram_equal_to_a_triangle_to_a_line}
\par\medskip\hrule
\end{document}
